%% kapa: exam=1 round=26 year=2015 period=1 form=A subject=civil_law subject_ko=민법 seq=10 number=10
%% source: Q-Net 감정평가사 기출문제 (cst00309) artlSeq=5138001 file=extracted/1008/5138001/제26회_감정평가사 1차_기출/제26회_감정평가사 1차_기출/1교시_감평_A형.pdf
%% conversion: kordoc 4.24.51 -> tools/gen1.py (auto); review=proofread-sonnet; flags=none
\begin{question}{10}
사회질서에 반하는 법률행위에 관한 설명으로 옳지 않은 것은? (다툼이 있으면 판례에 따름)
\begin{choices}
  \item 어떠한 일이 있어도 이혼하지 않겠다는 약속은 무효이다.
  \item 법정에 나와 증언할 것을 조건으로 대가를 지급하기로 약정한 경우, 그 대가의 내용이 통상적으로 용인될 수 있는 수준을 초과하면 그 약정은 무효가 된다.
  \item 이중매매임을 알고 부동산을 매수한 것만으로 제2 매매가 사회질서에 반하여 무효인 것은 아니다.
  \item 양도소득세의 일부를 회피할 목적으로 매매계약서에 실제로 거래한 가액보다 낮은 금액을 매매대금으로 기재한 경우에 그 매매계약은 무효이다.
  \item 반사회질서의 법률행위는 당사자가 그 무효임을 알고 추인하여도 새로운 법률행위를 한 효과가 생길 수 없다.
\end{choices}
\end{question}
\answer{4}
